\chapter{Misbehavior, Crises, Regulation and Self Regulation}

This lecture picks up exactly where the discussion of behavioral finance leaves us. If markets are shaped by human foibles, then we have to ask what institutions do about them. Regulation is the answer, but not in only one sense. It is directed, first, at manipulation, abuse, and the exploitation of human weakness; it is directed, second, at the more technical problem of keeping the financial system working at all. The lecture unfolds from that starting point: first \emph{too big to fail}, then the split between microprudential and macroprudential regulation, then the need for referees, and only after that the climb from regulation within the firm to regulation at the international level.

\section{Regulation After Behavioral Finance}

We begin, then, with the lecturer's explicit transition from the previous lecture. Regulation is not an afterthought to behavioral finance. It is what we do once we admit that finance is inhabited by imperfect people, by temptation, by persuasion, by manipulation, and by the possibility of plain bad character. At the same time, regulation has a more technical side: it deals with monopolies, externalities, and the need to prevent the whole financial system from malfunctioning.

The lecture chooses one system-level example almost immediately, because it gives the discussion urgency before any taxonomy appears. That example is \emph{too big to fail}. When a large financial firm collapses, the government often rescues it, not because anyone intended to create such a privilege, but because the failure of a sufficiently large firm can drag the whole system down with it. In a compact reconstruction,
\begin{align}
P(\text{bailout}\mid \text{large-firm failure})
>
P(\text{bailout}\mid \text{small-firm failure}).
\end{align}

That inequality is not a formal theorem from the board; it is a way of writing the lecture's causal claim. Once it is true in practice, large firms acquire an implicit government guarantee. The mechanism the lecture stresses is then
\begin{align}
\text{large-firm failure}
&\Rightarrow
\text{threat to the system}
\Rightarrow
\text{likely bailout},\\
\text{likely bailout}
&\Rightarrow
\text{implicit guarantee}
\Rightarrow
\text{greater risk-taking}
\Rightarrow
\text{greater systemic risk}.
\end{align}

That is the key analytical move. A rescue may stabilize one crisis in the short run, but the expectation of rescue changes behavior in advance. The guarantee is effectively free insurance provided by taxpayers, and it pushes risk upward. In that sense regulation is not only about stopping swindlers; it is also about stopping the system from training its largest institutions to gamble more aggressively.

\section{Too Big to Fail, Microprudential, and Macroprudential Regulation}

From there the lecture introduces a pair of terms that became especially prominent after the financial crisis of the 2000s: \emph{microprudential} and \emph{macroprudential}. ``Prudential'' here means regulation concerned with prudent standards of business conduct. The difference lies in the object of protection:
\begin{align}
\text{Microprudential regulation}
&\longrightarrow
\text{protect firms, households, and individual investors},\\
\text{Macroprudential regulation}
&\longrightarrow
\text{protect the system as a whole}.
\end{align}

The lecture is careful about this distinction. Microprudential regulation is the kind that protects small people, small firms, ordinary borrowers, and ordinary investors from being fooled or exploited. Macroprudential regulation is the kind that worries about crises that spread through the system. On this classification, \emph{too big to fail} is the paradigm macroprudential problem.

The lecturer also gives us a historical claim. Regulation has long involved both kinds, but until recently the mix leaned more heavily toward the microprudential side. Only after world financial crisis did the macroprudential side move to the center of the discussion. That historical change matters for the whole rest of the chapter, because many of the agencies at the end of the lecture are explicitly designed to think at the system level.

\subsection{Question \& Answer}

\paragraph{Question.} Why can rescuing a large firm create more risk rather than less?

\paragraph{Answer.} Because the rescue does not only affect the moment of collapse; it affects incentives before collapse. If the market learns that sufficiently large firms will be rescued to prevent system-wide failure, then those firms behave as though they possess a contingent public guarantee. They can therefore take risks that smaller firms cannot take. The immediate rescue may avert panic, but the expectation of rescue breeds the next round of excessive risk-taking.

\section{Why Markets Ask for Referees}

Before the lecturer turns to institutions, he pauses to justify regulation itself. He does not assume that the case for rules is obvious. Instead he says that regulators are like referees at a sports event. They enforce rules, decide when the rules have been broken, and punish violations. Players may complain about referees, but they still want them, because without them dangerous play would become nearly necessary in order to win.

That same idea is carried over to finance. In a competitive system, if everyone else is doing something shady, then one may feel compelled to do it too in order to survive. The lecture calls this a race to the bottom. In schematic form,
\begin{align}
\text{if others cheat}
\Rightarrow
\text{pressure to imitate}
\Rightarrow
\text{demand for enforceable rules}.
\end{align}

This is one of the lecture's most important normative points. Regulation need not be viewed as something imposed from outside against the wishes of all market participants. Competitors themselves may ask for it, because a well-policed game is preferable to a game in which success requires fouling.

The lecturer pushes the analogy further. Children on an empty field also create rules, and in effect referee one another. That is already a kind of self-regulation. He also insists that we should respect regulators as professionals. In sports we admire players more than referees; in finance we admire dealmakers more than regulators. But refereeing is a real craft, and regulation is a serious career. Some people move from Wall Street into regulation not because they failed, but because they would rather enforce decent rules than spend their whole lives in a morally compromised environment.

\subsection{Question \& Answer}

\paragraph{Question.} If competition is supposed to discipline behavior, why do competitors themselves ask for regulation?

\paragraph{Answer.} Because competition disciplines some margins and corrupts others. It may pressure firms to perform, but it can also pressure them to copy profitable misconduct. When rule-breaking becomes a competitive advantage, honest behavior is no longer individually sustainable. A referee solves that coordination problem by making restraint universal rather than optional.

\section{Five Levels of Regulation}

Only after those motivational steps does the lecture reset itself with a blackboard outline. On the left are the prudential terms; on the right is a hierarchy of regulatory scale. The board is not a formal diagram. It is an outline of the lecture's architecture.

\begin{figure}[t]
\centering
\includegraphics[width=0.88\textwidth]{lecture_12_figure_02.png}
\caption{Blackboard outline of prudential regulation and the five levels of regulatory scope.}
\end{figure}

The lecturer partially blocks the two prudential words on the left board, but the transcript makes the intended pair clear. A clean typographic reconstruction is
\begin{align}
\text{Too big to fail}
\qquad
\text{Microprudential}
\qquad
\text{Macroprudential}.
\end{align}

Beside that vocabulary the lecturer writes the five levels through which he will now climb:
\begin{enumerate}
\item Within firm
\item Trade groups
\item Local government
\item National
\item International
\end{enumerate}

It is important not to overread the board. There is no drawn arrow linking the prudential terms on the left to the five-part ladder on the right. The relation is conceptual rather than graphical. The prudential distinction tells us what regulation may be trying to protect; the five-level ladder tells us where regulation is located and how broad its scope has become.

The lecturer immediately adds a historical observation: the older forms of regulation are the narrower ones, and over time regulation has become broader in geographical scope. That is why international regulation is saved for the end. The lecture is not only a taxonomy; it is also a historical ascent.

\section{Within-Firm Regulation: Boards, Character, and Tunneling}

The first and lowest level is within the firm. The lecturer insists on counting this as regulation even though the government is not yet involved. Firms set up their own rules, and in the modern corporation the board of directors can be understood as a kind of internal regulatory body. He even says that boards could be called boards of regulators, though no one actually speaks that way.

The key element is the outside director. Inside directors work for the firm; outside directors do not, except through their board service. They are therefore the ones who make the board look most like a regulator. The lecture then ties this to the previous discussion of character. Some people have more character than others, and society tries, in its imperfect way, to keep people of weak character out of positions where they can do the most damage. A board of directors is part of that filtering mechanism. We place on it people with reputations for being serious, high-minded, and publicly accountable.

That is why the lecture pauses over the Yale Corporation. The point is not celebrity for its own sake. It is reputational merger. Richard Levin, Dan Malloy, Fareed Zakaria, Indra Nooyi, Mimi Gardner Gates: names like these signal that one is joining one's reputation to the reputation of the institution. Indra Nooyi also introduces the idea of interlocking directorship, because a director is judged partly by the other boards on which she is willing to serve. The lecture's point is that one would hesitate to propose a shady deal in front of such people.

The deeper reason for this internal regulation is that the corporation is extremely vulnerable to insider abuse. Once we move from a nonprofit example to a for-profit firm, the lecturer says, there are simply too many opportunities for sleazy behavior. That is why he introduces what he treats as perhaps the central internal pathology: \emph{tunneling}. Tunneling means taking value that belongs to the corporation and sneaking it into private pockets instead of leaving it with the shareholders who supplied the capital.

The internal logic can be summarized as
\begin{align}
\text{outside capital}
\Rightarrow
\text{delegated control}
\Rightarrow
\text{opportunities for tunneling}
\Rightarrow
\text{need for internal regulation}.
\end{align}

The lecture then accumulates examples, and the accumulation matters:
\begin{itemize}
\item \textbf{Asset sales.} The firm sells an asset below market value to a relative or associate, with compensation to the insider coming later and privately.
\item \textbf{Contracts.} The firm overpays a favored supplier, again with some return favor arranged off the books.
\item \textbf{Executive compensation.} The chief executive pays friends or cronies far too much and defends the arrangement as a legitimate judgment about merit.
\item \textbf{Expropriation of corporate opportunities.} A business opportunity that properly belongs to the firm is diverted into a new vehicle privately controlled by insiders.
\item \textbf{Insider trading.} Management uses nonpublic information, acquired because of its position in the firm, to trade for private gain.
\end{itemize}

The lecture wants us to feel how fragile the corporation would be if these channels were left uncontrolled. Shareholders do not monitor every asset sale, every contract, every compensation package, every deal opportunity, every piece of private information. A corporation is therefore, as the lecturer puts it, a delicate thing. It works only if there are institutions inside it that restrain tunneling.

\subsection{Question \& Answer}

\paragraph{Question.} Why would anyone invest in a corporation if insiders have so many ways to tunnel value out of it?

\paragraph{Answer.} Because the corporation is not just a productive unit; it is a governance arrangement. Boards of directors, especially outside directors, serve as internal regulators. They bring reputation, judgment, and community standards into the firm. The whole point is to make outside investors believe that managers cannot quietly convert corporate value into private value at every margin.

The lecture then sharpens the board's role by naming two legal standards:
\begin{align}
\text{Duty of care}
&:\ \text{act as a reasonable, prudent, rational person would act},\\
\text{Duty of loyalty}
&:\ \text{place the shareholders' interest ahead of private insider gain}.
\end{align}

The \emph{duty of care} means that the board member must get information, watch management, and exercise judgment. A board member is not the manager of the company, but neither is he a ceremonial guest who attends four meetings a year and nods at the chief executive's presentation. He must know enough to act as a regulator. The \emph{duty of loyalty} is usually interpreted, in the lecture's baseline formulation, as loyalty to shareholders. The lecturer notes that corporate social responsibility is widening this discussion, but he returns to the basic point: because the temptation to tunnel is so persistent, the board's first loyalty remains to those who put up the money.

That is why the lecturer treats within-firm regulation as perhaps the most essential form of regulation. Even before we move outward to exchanges or governments, the corporation itself must build mechanisms that make outside investment possible. Firms also create compliance departments, statements of purpose, and internal rulebooks. But the board remains the centerpiece, because it combines oversight, reputation, and fiduciary obligation in a single institution.

\section{Trade Groups, Local Regulation, and National Disclosure}

The lecturer then says explicitly that he wants to move up another step. We leave the single firm and turn to trade groups: organizations formed by firms themselves, sometimes as self-regulatory bodies. The historical example he chooses is the New York Stock Exchange, founded in 1792 after the first American stock market crash.

The Buttonwood Agreement matters here because the lecturer actually reads it and forces us to confront its blunt economic content. It is not a long ethical manifesto. It is a very short agreement that, in essence, establishes a minimum commission and gives preference to negotiation among the members themselves. The commission term can be summarized as
\begin{equation}
c_{\text{Buttonwood}}=\frac{1}{4}\%=0.25\%=0.0025.
\end{equation}

Taken narrowly, that looks like price-fixing combined with exclusion of outsiders. In other words, it looks like a cartel. But the lecture does not leave the matter there. The more charitable interpretation is that the stockbrokers were reacting to a speculative boom and crash associated with a dubious promoter, and they were trying to exclude bad actors while maintaining incomes high enough to support what they regarded as responsible conduct. The lecturer is explicit that this second interpretation is partly reconstructed between the lines. The tension is the point: a trade group can stabilize conduct and create monopoly power at the same time.

\subsection{Question \& Answer}

\paragraph{Question.} How can self-regulation help markets without simply turning into cartelization?

\paragraph{Answer.} Because the two functions are genuinely separable even when they arrive in the same institution. A trade group can screen participants, police conduct, and build reputational norms. But it can also fix prices and exclude rivals. Public regulation then has to do two things at once: preserve the useful parts of self-regulation, such as surveillance and ethical rule enforcement, while dismantling monopoly privileges.

That later dismantling comes with \emph{May Day}, May 1, 1975, when fixed commissions on the New York Stock Exchange were made illegal. The lecturer's point is not that the exchange disappeared. It is that the old cartel structure was broken. Members could no longer maintain mandatory commissions and closed dealing relations among themselves. The national market system widened the broker's obligation: not loyalty to a closed club, but duty to seek the best price for the customer on any exchange. The British \emph{Big Bang} of 1986 plays the same role in the lecture's comparative story. In both cases, deregulation of commissions is itself a form of regulation against cartelization.

What remains even after that change is surveillance. The New York Stock Exchange still monitors trading records for manipulation. So the lecturer does not treat the post-1975 world as a world without referees. It is a world in which one particular kind of self-regulation---fixed commissions and monopoly dealing---has been rejected, while surveillance and conduct policing remain.

From trade groups the lecture moves to local regulation. In the United States, before the 1930s, regulation was overwhelmingly local. The federal government did very little. States developed \emph{Blue Sky Laws} during the Progressive Era, beginning in Kansas in 1911, to block fraudulent or wildly overpromoted securities. The colorful name is tied, at least in one common explanation, to sales pitches promising that prices could rise with nothing above them but the blue sky.

Local regulation, however, ran into a structural problem:
\begin{align}
\text{state-by-state regulation}
\Rightarrow
\text{cross-state evasion}
\Rightarrow
\text{need for national regulation}.
\end{align}

That is the mechanism the lecture wants us to see. Once telephones and boiler rooms make it easy to sell worthless securities across state lines, fragmented state enforcement becomes inadequate. A boiler room in one state can pressure victims in another state, and the geographical fragmentation of law becomes an invitation to fraud. So the Blue Sky system remains historically important, but it is no longer sufficient.

The next step is national regulation under the New Deal. In 1934 the United States created the Securities and Exchange Commission. The lecturer does not present it merely as another agency on a list. He makes it vivid through figures such as William O. Douglas and Arthur Levitt, each of whom, in his own way, confronted Wall Street's hostility to disclosure-based regulation. The guiding principle comes from Louis Brandeis: ``sunshine is the best disinfectant.'' In finance that means the following:
\begin{align}
\text{secrecy}
\Rightarrow
\text{scope for abuse},
\qquad
\text{disclosure}
\Rightarrow
\text{public scrutiny}
\Rightarrow
\text{cleaner markets}.
\end{align}

That is why the SEC is built around disclosure. Public companies must file documents; those documents are made available through EDGAR; and, in principle, anyone can inspect them. The lecture returns to this point because it wants disclosure to be understood not as bureaucratic clutter but as a regulatory technology.

The SEC also draws the line between public and private securities. A firm that wants to go public must pass through the regulated procedure of an IPO. A private fund is given more freedom, but only on the condition that access is restricted. The lecture describes two broad hedge-fund categories in approximate quantitative form:
\begin{align}
N_{\text{class 1}} &\le 99,
&
y_{\text{single}} &\ge 200{,}000,
&
y_{\text{married}} &\ge 300{,}000,
&
W_{\text{investable}} &\ge 1{,}000{,}000,
\\
N_{\text{class 2}} &\le 500,
&
W_{\text{individual}} &\ge 5{,}000{,}000,
&
W_{\text{institution}} &\ge 25{,}000{,}000.
\end{align}

The statutory labels are noisy in the transcript, so the important point is conceptual: private funds are permitted greater latitude because participation is restricted to wealthy investors presumed able to protect themselves.

The lecture then uses the fee formula ``2 and 20'' to make that world feel concrete. If \(A\) denotes assets under management and \(\Pi\) denotes trading profits, a cautious reconstruction is
\begin{equation}
f(A,\Pi)=0.02A+0.20\max\{\Pi,0\}.
\end{equation}

\paragraph{Worked example.} If \(A=10^9\), then the management fee alone is
\begin{equation}
0.02\times 10^9 = 2\times 10^7 = 20{,}000{,}000.
\end{equation}
That is the lecturer's arithmetic check. Even before performance fees, the manager's compensation can be enormous. Losses do not come symmetrically out of the manager's pocket. The point is not to build a theory of hedge-fund contracting. The point is to dramatize the two-tier system of investor protection: one heavily regulated public tier for ordinary investors, and one much looser private tier for wealthy investors.

The national section ends, appropriately, with surveillance. The SEC does not only require disclosure; it also watches the market, often together with self-regulatory organizations. The lecture's Lotus case is a compact illustration:
\begin{equation}
N_{\text{Lotus tippees}}\approx 25,
\qquad
V_{\text{Lotus trades}}\approx \$500{,}000.
\end{equation}
A private comment about IBM's plan to acquire Lotus spreads through a network of acquaintances, trading surges, phone records are subpoenaed, and the chain is reconstructed. The Emulex case makes the complementary point: false information can also be used to manipulate markets, and surveillance can trace the trade, the short cover, the computer, and the person who sent the fake press release. In both stories the lecture returns us to its earlier analogy. Markets do not function well without referees.

Two final institutions complete the national level. FASB, recognized by the SEC, defines GAAP, providing accounting standards for disclosure. SIPC, created in 1970 after brokerage failure exposed the fragility of accounts held in street name, protects brokerage customers in a role parallel to that played by deposit insurance for bank depositors. Both are examples of how national regulation is partly delegated and partly specialized.

\section{From Dodd--Frank to Europe and the International Layer}

The final movement of the lecture shifts from the older national framework to the post-crisis expansion of macroprudential and international regulation. In the United States the key event is the Dodd--Frank Act of 2010. The lecturer reads it as a push away from a predominantly microprudential regime and toward a more macroprudential one. Earlier legislation had already addressed some fraud against individuals. Dodd--Frank adds institutions whose job is explicitly to worry about systemic risk.

The clearest example is the Financial Stability Oversight Council, FSOC. Its purpose is to think about \emph{too big to fail} and about risks that threaten the whole system rather than merely one individual contracting party. Alongside it stands the new consumer-protection bureau created by Dodd--Frank, whose purpose is more microprudential: to focus specifically on ordinary borrowers, mortgage customers, and credit users who were not well served by the older regulatory structure. The pair reproduces, at the institutional level, the distinction made near the start of the lecture.

The same broadening then appears in Europe. The European Supervisory Framework, also associated with 2010, includes several new bodies:
\begin{itemize}
\item the European Systemic Risk Board in Frankfurt;
\item the European Banking Authority in London;
\item the European Securities Market Authority in Paris;
\item the European Insurance and Occupational Pension Authority in Frankfurt.
\end{itemize}

The lecturer is careful here too. These institutions are new, and at the time of the lecture their eventual functioning is not yet fully known. Still, the direction is unmistakable: regulation is being lifted from the level of separate national jurisdictions toward a wider economic region.

From there the lecture goes fully international. The problem is easy to state. If regulation is costly, then financial firms can move offshore. If they dislike U.S.\ rules, there is always the Bahamas or the Cayman Islands. So we are led to another causal chain:
\begin{align}
\text{tighter national regulation}
\Rightarrow
\text{offshore migration}
\Rightarrow
\text{need for international coordination}.
\end{align}

This is why the international layer matters. The Bank for International Settlements, in Basel, does not legislate for the world, but it convenes and coordinates. The lecture says it brings together roughly \(57\) central banks. The Basel Committee operates similarly through recommended banking standards rather than direct command. Its major sequence is
\begin{equation}
\text{Basel I }(1988),\qquad
\text{Basel II }(2004),\qquad
\text{Basel III }(2010).
\end{equation}

The lecture's tone here is sober. International bodies do not usually possess the direct coercive force of national regulators. Their authority comes through recommendation, coordination, and repeated agreement. But that does not mean they are unimportant. After a global banking crisis, recommendation itself can have real force if the major countries agree to take it seriously.

A chronology helps keep the lecture's historical scale in view:
\begin{align}
1792 &: \text{Buttonwood Agreement and the New York Stock Exchange},\\
1911 &: \text{Kansas Blue Sky law},\\
1934 &: \text{Securities and Exchange Commission},\\
1970 &: \text{SIPC},\\
1975 &: \text{May Day and the end of fixed commissions},\\
1986 &: \text{Big Bang in the United Kingdom},\\
2008 &: \text{G20 expansion},\\
2010 &: \text{Dodd--Frank, European supervisory reforms, Basel III approval}.
\end{align}

The path from the G6 to the G7 and then to the G20 marks the same broadening in political form. Finance ministers first coordinate among a narrow group of large countries, then among a broader one. The G20, in the lecture, represents a more global forum, one that now includes countries such as China and India rather than overwhelmingly reflecting Europe and North America. The Financial Stability Board, also in Basel, is presented as the G20-linked body making recommendations for global regulatory procedures.

The lecture closes with guarded optimism. We are seeing financial regulation become larger in scale because finance itself has become larger in scale. A worldwide crisis has forced regulation to think not only about fraud against individuals, but about the stability of the world economy. That, for the lecturer, is an inspiring development, provided the international cooperation continues.

\section{Summary}

The lecture begins with human misbehavior and ends with world-scale institutions, and the movement between those two endpoints is carefully staged. We start with the thought that regulation must answer both to psychology and to system design. We then see why \emph{too big to fail} creates moral hazard and why that problem is macroprudential. We pause to ask why markets want referees at all, and the answer is that competition without enforceable rules can become a race to the bottom.

Only then do we climb the five levels of regulation. Within the firm, boards of directors, reputation, and fiduciary duties hold back tunneling. At the level of trade groups, self-regulation can improve conduct but also harden into cartelization. At the local level, Blue Sky laws matter historically but fail once finance easily crosses state lines. At the national level, the SEC builds regulation around disclosure, surveillance, and investor protection. And at the international level, the lesson of crisis is that finance outgrows any one jurisdiction, so regulation must become more coordinated, more macroprudential, and more global.

That is the real shape of the lecture. It is not a miscellaneous list of agencies. It is a guided ascent from the moral fragility of the individual firm to the systemic fragility of the world financial order.